\subsection{Two-mark angular algebra}\label{c4:pair-angular}

We next express the numerator of \eqref{c4:Pair-sum} as a radial series.
The two color blocks require two coupled orbits of charges.  We give the
charge space, the field phases, and the conjugation explicitly; these data
also determine the partially suppressed trace used later.

\paragraph{The charge space and its field phases.}
The space has basis symbols \(\alpha_{ij}=C^j\alpha_{i0}\) (\(i=1,2\), denoting sectors \(R,L\), and \(0\le j<8\)), quotiented by the radical. The two diagonal blocks of the symmetric Gram are circulant with first row \((2,-1,1,0,0,0,1,-1)\), cross block \((0,0,0,-2,2,-2,0,0)\); \(C\) is the simultaneous forward shift. Write \(\mathcal L\) for the integral span as usual. Define \(h_1\) by pairings \((1,0,0,0,0,0,1,0)\) in its sector, \((0,0,0,-2,0,0,0,0)\) in the other, \(h_2\) symmetrically; \(v\) by \((0,0,0,-1,0,0,1,0)\) in both. Let \(J\) be the alternating form with within-sector first row \((0,1,1,0,0,0,-1,-1)\) in sector 1, its negative in sector 2, zero across. Put \(p_C=\sum C^j h_1/8\). These descend as indicated: the nullspace consists of span of sector-1-even plus sector-2-odd sums and the reverse sum (diagonalize by eighth roots); \(p_C^2=1/4,\ (p_C,\alpha_{ij})=(-1)^{i+1}/4\). The \(J\)-pairing and ordinary pairing coincide modulo 2 on \(\mathcal L\times(\mathcal L+\sum_i\mathbb Z h_i)\); and \(a^t J v+(a,v)+3(a,p_C)\in2\mathbb Z\) there for \(a\in\mathcal L\). All these checks just multiply the circulant data (in solving for pairings remove the two null vectors).

Use the same Fock states and \(W_b\) as before on \(\mathcal L^*\), but with momentum phase
\[
 \exp(\pi i[b^t J P/2-3k_b(p_C,P)])
\]
(\(P\) momentum, \(k_b=0\) for lattice and color charges, \(k=1\) per mark \(v\), \(-2\) for compensator, extended additively and preserved by \(C\)). The oscillator trace is twisted initially just by \(\widehat C Q_L^H\) (no other sign), scaled as before; for stationary trace comparisons we also allow a character \(e^{2\pi i(\sigma,P)}\) with \(\sigma=0\) or \(2v\). Use \(X_b=[z^0]W_b\) for lattice roots of norm square 2. The same residue commutators hold with the new cocycle (coefficient \(e^{\pi i a^t J b/2}\) for root \(a\) acting when \((a,b)=-1\)), by the parity identities above.

\paragraph{The finite conjugation.} Within each sector \(i\) introduce labels \(j=1,2,3,4,E\) with
\[\begin{gathered}
 r=(\alpha_{i0},0,-\alpha_{i1},\alpha_{i2},-\alpha_{i1}+\alpha_{i7}),\\
 e_{jk}=r_j-r_k,\quad H_k=h_i+r_k-r_1,\quad E_{jk}=e^{-\pi i e_{jk}^t J H_k/2} X_{e_{jk}} .
\end{gathered}\]
Use this only for \(1234\) and the extra \(E3,E4\), which are roots. The \(1234\) differences have ordinary \(A_3\) root pairings, with \((e_{jk},H_l)=\delta_{jl}-\delta_{kl}\); across sectors these two \(A_3\)'s are orthogonal. Thus we again get commuting \(SL_4\) operators with matrix action on the four fields \(W_{H_l}\). Also \([E_{E3},E_{34}]=E_{E4}\), central for those three, by the residue formula. The exponentiation identities here work as before, using finite-dimensional modules along positive root subspaces (including the \(A_2\) form on the \(E3,34\) span). In particular no \(SL_5\) claim on the entire list including \(E\) is needed.

Set \(c=q^{-3}\). In each \(SL_4\) use lifted operators \(G_{0i},\widehat S_i,D_i\) for
\[
\begin{gathered}
 G_0=\begin{pmatrix}1&0&0&0\\c&1&0&0\\c^2&c&1&0\\0&0&0&1\end{pmatrix},
 \qquad S=\begin{pmatrix}0&-1\\1&0\end{pmatrix}\text{ on }12,\\
 D_i=\exp(-\pi i(\alpha_{i0},P)/4)\quad(\operatorname{diag}(-1/c,-c,1,1)).
\end{gathered}
\]
Put \(t=-c^2,\ k_1=-\mathrm{i}/c,\ k_2=\mathrm{i}/c,\ V_i=\exp(k_iE_{E3})\exp(t E_{34})\), \(g=G_{01}G_{02}V_1 V_2\). These two sector actions in \(g\) commute by nonnegative cross pairings (and orthogonality for \(A_3\) across sectors). Then \(g\) sends \(W_{h_i}\) by conjugation to the desired \(W_{H_1}+cW_{H_2}+c^2W_{H_3}\), leaving \(W_v,W_{-2v}\) fixed; indeed the \(V_i\) commute with these initial pure fields and \(G_0\) uses just the three-label subalgebra. Write \(A=S_{\alpha_{10}}S_{\alpha_{20}}C,\ \widehat A=\widehat S_1\widehat S_2\widehat C\).

We have
\[
 g^{-1}\widehat Cg=\mathcal U D_1D_2\widehat A
\]
where the unipotent acts by a product of \(\exp(u X_b)\) (unit-modulus or fixed complex coefficients immaterial except as specified below), with rightmost-first root order across sectors grouped as: \(b=A e_{34}\) (call it \(e_{1F}\)), then \(e_{12}, e_{24}\), then \(e_{E3},e_{E4}\), each in both sectors. This fact is used with separate small contour perturbations within the groups. Further, if suppressing just \(1F,12,E3\) in sector 1, the resulting entire twist equals
\begin{equation}
 \exp(-\mathrm{i}c E_{E4}^{(1)})\exp(c E_{24}^{(1)})\,
     g_2^{-1}D_1\widehat S_1\widehat C g_2,\qquad g_2=G_{02}V_2,                    \label{c4:Pair-Bcut}
\end{equation}
where the initial two factors also commute with \(g_2\).

We verify the factorization, including the surviving coefficients in
\eqref{c4:Pair-Bcut}.  Their values will matter in the partial trace. In either sector \(\widehat C\) sends the \(123\) algebra by conjugation to negative transpose with labels sent to \(324\), conjugated by \(\operatorname{diag}(-1,1,-1)\) on the renamed triple (indeed \(E_{12}\to E_{23}, E_{23}\to E_{42}\)); it sends \(E_{E3}\) to \(s_i E_{12}\), \(s_1=\mathrm{i},s_2=-\mathrm{i}\). These follow directly using the normalization: phase ratios for those three generators in units \(\pi i\) are \(0,0,(-1)^{i+1}/2\). Thus \(\widehat A\) sends \(E_{E3}\to-s_i E_{21}\). Also \(\widehat A\) sends \(E_{34}\) to a scalar root operator of charge \(e_{1F}\) (justified below). Now \(V^{-1}\prod_i U_{0i}(\widehat A V\widehat A^{-1})\), where \(V=V_1V_2\) and \(U_0=G_0^{-1}(\widehat C G_0\widehat C^{-1}) S^{-1}\) in each group, multiplies \(\widehat A\). Rearrange on the left first moving the \(E3\)'s outward: this gives \(\exp(-k_i E_{E3}-t k_i E_{E4})\) per sector on the far left, leaving \(\exp(-t E_{34})\). On the far right before \(\widehat A\) retain both \(1F\) root exponentials, leaving inner \(\exp(-E_{21}/c)\)'s. This ordering uses cross-commutation of the generating \(V\)-terms. In each group the middle computation now is
\[
 (1-t E_{34})G_0^{-1}(\widehat C G_0\widehat C^{-1}) S^{-1}
 (1-E_{21}/c)=(1+cE_{24})(1-E_{12}/c)D_i
\]
in ordinary matrix notation. Move \(D_1D_2\) to the outgoing \(\widehat A\), just scaling the \(1F\)'s. This proves the order and coefficients. Keeping only \(E4,24\) in sector 1, move those surviving factors to the far left (nonnegative cross pairings); the calculation in sector 2 then reverses to \eqref{c4:Pair-Bcut}. Indeed \(D_1\) commutes with the inner sector-2 algebra, \(g_2\) and the \(E3,E4\) row factors there; the \(D_1\)-conjugation of sector-2 \(1F\) is retained as prescribed. Likewise any indicated character commutes with the conjugations by integrality.

\paragraph{The Weyl phases.} For the Weyl operator \(\widehat S_i\) the action reflects oscillator modes and sends a pure momentum \(n\to S_\alpha n\) with phase
\[
 \exp(\pi i[k\,\alpha^t J(h_i-n)+k^2-k]/2),
 \qquad \alpha=\alpha_{i0},\quad k=(\alpha,n).
\]
For \(k\ge0\) use \(E_{21}^k/k!\) on the highest state, as in the previous angular calculation; the Vandermonde contributes \((-1)^{\binom{k}{2}}\). For \(k<0\) use \((-E_{12})^{-k}/(-k)!\); the same formula results. Modes transform by their adjoint triples as before. This shows conjugation carries lattice root operators to scalar multiples of reflected ones: the possible extra momentum dependence in the phase relative to the new cocycle reduces to \(k[(\alpha,b)-\alpha^t J b]\in2\mathbb Z\). It also shows the total \(D_1D_2\widehat A\) uses oscillator action \(A\) and just the product of these two Weyl phases on \(Cm\) for starting momentum \(m\), and the displayed diagonals on \(A m\). Finally \(C h_i=A h_i,\ C v=A v\), and these images pair nonnegatively with all indicated unipotent roots, orthogonally with the Weyl roots. Thus the fields for a shifted color or mark group, conjugated by \(g^{-1}\widehat Cg\), are exactly the pure images under \(C\), without further scalar.

